% ─────────────────────────────────────────────────────────────────────────────
% Environment: Cyberlab BFH (real network, ~35 live hosts on 193.5.80.0/24)
% claude-opus-4-7 (multi-agent + unstructured), qwen3:8b (local) + static reference
% Evaluation criteria: see §\ref{subsec:criteria}
% ─────────────────────────────────────────────────────────────────────────────
\section{Network Security Lab BFH}\label{sec:eval-cyberlab}

\paragraph{Overview}


\begin{table}[H]
    \caption{Benchmark Results --- Network Security Lab BFH}
    \centering
    \resizebox{\textwidth}{!}{
        \begin{tabular}{l c c c c c c}
            \hline\hline
            Configuration & Avg.\ Duration & Avg.\ Tokens & Avg.\ Hosts & Avg.\ Services & Avg.\ Findings & Success Factor \\ [0.5ex]
            \hline
            Multi-Agent (claude-opus-4-7)  & 510 s  & 116,738  & 28 & 53 & 52 & 5/27   \\
            Unstructured (claude-opus-4-7) & 739 s  & 222,094 & -- & -- & -- & 6/109  \\
            Multi-Agent (qwen3:8b, local)  & ---    & ---     & -- & -- & -- & 0/6    \\
            Static Scenario                & 1780 s & ---     & -- & -- & -- & --     \\
            \hline
        \end{tabular}
    }
    \label{tab:overview-nslab}
\end{table}

Table~\ref{tab:overview-nslab} summarises the quantitative benchmark values,  the configurations are discussed below.


\subsection{Quantitative Criteria}\label{subsec:quantitative-criteria-cyberlab}

\begin{figure}[H]
    \centering
    \includesvg[width=0.7\linewidth]{figures/cyber-lab/benchmark_duration_comparison_ns_lab.svg}
    \caption{Mean run duration per model (Cyberlab BFH).}
    \label{fig:benchmark_duration_cyberlab}
\end{figure}

Figure~\ref{fig:benchmark_duration_cyberlab} shows the mean run duration per configuration.
The multi-agent setup completes in 510\,s on average, while the unstructured output scenario
requires 739\,s — more than twice as long.
The extended duration of the unstructured version is due to the fact from its tendency to produce lengthy
free-form responses before each tool invocation.
The static benchmark scenario needs 1780\,s to enumerate the whole network which exceeds the \gls{AI} configurations by far.


\begin{figure}[H]
    \centering
    \includesvg[width=0.7\linewidth]{figures/cyber-lab/benchmark_token_usage_ns_lab.svg}
    \caption{Mean total token usage per model (Cyberlab BFH).}
    \label{fig:benchmark_tokens_cyberlab}
\end{figure}

Figure~\ref{fig:benchmark_tokens_cyberlab} visualises mean token consumption.
The unstructured output scenario consumes 222,094 tokens on average, nearly three times the
116,738 tokens used by the multi-agent configuration, roughly 100\% higher token consumption.

\begin{figure}[H]
    \centering
    \includesvg[width=0.7\linewidth]{figures/cyber-lab/benchmark_tokens_vs_duration_ns_lab.svg}
    \caption{Total tokens vs.\ run duration per individual run (Cyberlab BFH).}
    \label{fig:benchmark_scatter_cyberlab}
\end{figure}

Figure~\ref{fig:benchmark_scatter_cyberlab} plots total token consumption against run duration
for every run.
The unstructured output runs shows high variance, spanning 310\,s to
1{,}091\,s and 119\,k to 305\,k tokens.
Multi-agent runs cluster more tightly, with durations between 84\,s and 684\,s.
The positive correlation visible within each group confirms that longer runs consume more
tokens.

\begin{figure}[H]
    \centering
    \includesvg[width=0.7\linewidth]{figures/cyber-lab/benchmark_services_ns_lab.svg}
    \caption{Average hosts, services, and findings discovered per model (Cyberlab BFH).}
    \label{fig:benchmark_services_cyberlab}
\end{figure}

Figure~\ref{fig:benchmark_services_cyberlab} compares services, host and findings across the two
configurations.
The multi-agent setup discovers an average of 28 hosts, 53 services, and 52 findings per
run on the Cyberlab BFH segment.
The unstructured output scenario does not produce structured discovery metrics, as its
assessment is written as free-form text rather than a machine-readable result map, so no
comparable figures are available for that configuration.
The static scenario lists again every service per host which result in more services (130) but rather close finding count with
52 findings for \gls{AI} vs 58 for the static scenario.

\subsection{Qualitative Criteria}\label{subsec:qualitativ-criteria-cyberlab}

The evaluation criteria and scoring scale are defined in table~\ref{tab:qualitative-scale}.
The qualitative findings are not evaluated in detail here, because the BFH network security lab is still in use.
The most severe findings include the full extraction of end-of-life systems, open \gls{DNS} recursion, and cloned host keys.
Flaws in this network are partly introduced to train the students which further complicates the evaluation.

Three findings serve as quality indicators:

\begin{table}[h]
\centering
\caption{Qualitative Evaluation -- Cyberlab BFH}
\label{tab:qualitative-results-cyberlab-bfh}
\small
\begin{tabular}{@{}lccl@{}}
\toprule
\textbf{Model} & \textbf{Correctness} & \textbf{Hallucination} & \textbf{Notes} \\
\midrule
claude-opus-4-7 (multi-agent)       & 7 & 2 & All major structural risks found; SNMP, open DNS recursion and cloned keys missed \\
claude-opus-4-7 (unstructured)      & 9 & 2 & All three discriminating findings surfaced; ESXi host discovered \\
qwen3:8b (local, multi-agent)       & 1 & N/A & Zero CLI calls; no reconnaissance executed \\
Static Scenario (reference)         & 10 & N/A & Deterministic script; exhaustive enumeration; no assessment text \\
\bottomrule
\end{tabular}
\end{table}

\todo{the table is missplaced}

\textbf{claude-opus-4-7 in multi-agent mode} achieves a correctness score of 7.
In the five successful runs the model consistently identifies the six high-severity structural
risks present in the network.
All findings are accompanied by severity ratings and actionable remediation steps.
However, none of the successful runs includes \gls{UDP} enumeration, so all flaws on this protocol went undiscovered.
Similarly, \mintinline{bash}{dns-recursion} and \mintinline{bash}{ssh-hostkey} comparison scripts are absent from
most runs, leaving the open resolver risk and the cloned key finding unexplored.
The hallucination score is 2: every claim in the final reports is directly traceable to
\mintinline{bash}{nmap} \gls{NSE} output, with no fabricated CVE identifiers or version numbers.
A notable reliability concern is the high error rate of the multi-agent setup: only 5 of 27 attempts produced a fully structured result.

\textbf{claude-opus-4-7 in unstructured output mode} achieves the highest correctness score of
9 among the \gls{AI} configurations.
When successful (six of 109 attempts), the model issued \gls{UDP} scans
alongside \gls{TCP} and discovered all the flaws listed above.
The hallucination score is equally 2: all assessment statements are grounded in concrete scan
artifacts.
The main limitation of this configuration is its very low success rate (5.5\,\%) and the
substantially higher token cost (mean 222,094 tokens), driven by the unstructured reasoning
traces produced before each tool call.

\textbf{Qwen3:8b} is not further assessed due to its low success rate, which is nearly zero.
